% Laget automatisk av elevark/lag_elevark.py fra sim/tog-og-tunnel.html. Endringer her blir overskrevet.
\documentclass[11pt]{article}
\usepackage{elevark}
\begin{document}
\elevarkhode{Fysikk 2 · Spesiell relativitetsteori}{Tog og tunnel}{https://lektorbo.no/sim/tog-og-tunnel.html}{Simuleringen}

Et tog er 200~m langt og en tunnel er 150~m lang, målt i ro. Toget kjører gjennom tunnelen med en fart nær lysfarten. I tunnelens system er toget kontrahert, og hvis det er kort nok, kan begge portene lukkes samtidig mens hele toget er inne. Men i togets system er det tunnelen som er kontrahert. Hvordan kan begge ha rett?

\section*{Oppgaver}

Start med farten $0{,}80c$ i tunnelens system.

\begin{oppgaveliste}
\oppgave \fase{test}{Regn}Regn ut $γ$ og lengden av toget i tunnelens system. Passer toget inn i tunnelen? Kjør simuleringen og sjekk.
\svar
\oppgave \fase{gjett}{Gjett}Bytt til togets system. Hvor lang er tunnelen nå? Kan toget være helt inne i tunnelen? Hva tror du skjer med portene?
\svar
\oppgave \fase{forklar}{Forklar}Se på tabellen over hendelser. I hvilken rekkefølge lukkes og åpnes portene i togets system? Forklar hvorfor toget ikke blir truffet av noen port i noen av systemene.
\svar
\oppgave \fase{test}{Mål}Velg «Tunnelfolket måler toget». Hvor langt er det mellom merkene på bakken? Bytt til togets system. Hvorfor mener togfolket at målingen er feil? Se på tidene for merkene.
\svar
\oppgave \fase{test}{Mål}Velg «Togfolket måler tunnelen». Hvor langt er det mellom merkene på toget? Hvem måler den andre som kortest?
\svar
\oppgave \fase{test}{Undersøk}Hvor stor må farten minst være for at toget skal passe inn i tunnelen sett fra tunnelen? Regn ut, og test.
\svar
\oppgave \fase{utfordring}{Utfordring}Sett farten til $0{,}60c$. Hva skjer nå i tunnelens system? Skjer det samme i togets system? Hvorfor må alle være enige om om porten treffer toget?
\svar
\end{oppgaveliste}

\end{document}
